\section{Relative Finite-Pattern Realization}
\label{sec:finite-separation}

The free-completion theorem gives canonical maps into compatible total
extensions, but it does not say which distinctions between suspended Answers
are visible in independently specified models.  We now restrict attention to
observers that preserve the base pointwise and have only finitely many states
outside it.

Let $T$ be a compatible total extension with injective base embedding
$i:A\hookrightarrow T$.  Its genuine outside part is
\[
 \operatorname{Out}_D(T)=\{z\in T\mid \forall a\in A,\ z\ne i(a)\}.
\]
An intrinsic observer has budget $n$ when
$\operatorname{Out}_D(T)$ injects into $\Fin(n+1)$.  The base $A$ may be
infinite.  The budget counts only states beyond the fixed copy of $A$.

A canonical stage-$n$ representative has carrier
\[
 A\sqcup\Fin(n)\sqcup\{\overflow\},
\]
with total operations that preserve every application defined in $D$.  The
formal representation theorem identifies separation by arbitrary intrinsic
budget-$n$ observers with separation by these finite-tag models.  Write
$x\stageeq{n}y$ when every stage-$n$ finite-tag model interprets $x$ and $y$
equally.

The original pairwise separator can be strengthened.  Instead of selecting the
subterms of two roots, select the union of the finite subterm closures of an
arbitrary finite family.  Old leaves keep their base values.  Every selected
suspended node receives the tag of its first occurrence in the selected list,
and all other behavior is sent to overflow.  A finite lookup table reproduces
each selected normal suspension exactly.  Induction over the selected closure
then shows that every selected generated Answer evaluates to its injective
encoding.

\begin{theorem}[Relative finite-pattern realization]
\label{thm:finite-pattern-realization}
Let $x_1,\ldots,x_k\in\RAlg{D}$ be any finite family.  There exist $n$ and one
compatible stage-$n$ finite-tag observer $T$ such that the canonical
interpretation
\[
 \Phi_T:\RAlg{D}\to T
\]
is injective on $\{x_1,\ldots,x_k\}$.
\end{theorem}

The construction fixes the entire base carrier pointwise; only the selected
suspended pattern consumes finite tags.

\begin{corollary}[Intrinsic finite-complement separation]
\label{thm:finite-external-separation}
For every distinct $x,y\in\RAlg{D}$ there is a compatible total extension
whose outside part has at most $|x|+|y|+1$ states and in which $x$ and $y$ have
different interpretations.  Equivalently,
\[
 \bigl(\forall n,\ x\stageeq{n}y\bigr)\Longleftrightarrow x=y.
\]
\end{corollary}

The finite-pattern theorem is structurally stronger than the pairwise bound,
while the pairwise implementation retains the sharper explicit estimate
$n=|x|+|y|$.

\begin{proposition}[Product obstruction over an infinite fixed base]
\label{prop:relative-product-obstruction}
The carrier-side condition ``pointwise-fixed base with finite outside
complement'' is not closed under ordinary binary products.  Concretely, each
embedding
\[
 \mathbb N\hookrightarrow \mathbb N\sqcup\{\star\}
\]
has a one-point outside complement, whereas the diagonal embedding of
$\mathbb N$ into
$(\mathbb N\sqcup\{\star\})^2$ has infinite outside complement.  Indeed, the
mixed points $(a,\star)$, for $a\in\mathbb N$, form an injective copy of
$\mathbb N$ outside the diagonal old image.
\end{proposition}

Thus the usual residual argument that combines finitely many pairwise
separators by taking their Cartesian product does not stay inside the
present observer class when the preserved base is infinite. The direct
finite-pattern construction in \cref{thm:finite-pattern-realization} is
therefore not replaced by that standard product trick.

\begin{remark}[What is and is not being claimed]
The selected-subterm plus sink construction has standard automata-theoretic
antecedents \cite{ComonEtAl2008}.  The contribution claimed here is not the
existence of a sink automaton in isolation.  It is the relative realization
statement for observers that must preserve a possibly infinite partial base
pointwise while budgeting only the complement.  In particular, the codomain
need not be finite when $A$ is infinite.
\end{remark}

If the base admits an injection into $\Fin(m)$, then the same observer becomes
an ordinary finite algebra with total-state bound
\[
 m+|x|+|y|+1.
\]
Hence $\RAlg{D}$ is residually finite in the ordinary sense for finitely coded
bases.  This qualitative finite-base consequence is classical in the
finite-signature ground-equational setting; the relative construction is used
because it remains meaningful when the base itself is infinite.

\begin{proposition}[Closure of the base image]
\label{prop:old-image-closed}
The image of $A$ in $\RAlg{D}$ is closed under every generated total operation
if and only if the evaluator of $D$ is total.
\end{proposition}

Thus in a genuinely partial example the preserved base is not a total
subalgebra of $\RAlg{D}$.  This is one reason ordinary finite-Rees-index
language does not literally describe the relative observer situation.
